\documentclass[11pt]{article}
\usepackage[margin=1in]{geometry}
\usepackage{amsmath}
\usepackage{siunitx}
\usepackage{hyperref}

\title{How Much Dark Matter Could Be Bound to Earth?}
\date{}
\begin{document}
\maketitle

The amount of dark matter (DM) gravitationally bound to Earth is not fixed by the local Galactic DM density alone. It depends on how efficiently particles are captured and retained, and on the spatial distribution of the bound population. A useful way to state the question is therefore to ask what halo masses remain compatible with a specified observational limit and a model for the bound states.

As a reference hypothesis, take the bound DM mass to be $0.38\%$ of Earth's mass:
\begin{equation}
  M_{\rm DM}=3.8\times10^{-3}M_\oplus
  \simeq 2.27\times10^{22}\,\mathrm{kg}.
\end{equation}
This is substantial on particle-physics scales, but small compared with Earth's mass. It is an assumed benchmark, not a value derived by the numerical code or established by the shell constraint. The distinction matters: the benchmark describes an integrated mass, whereas an experiment constrains DM in a particular region, and different radial profiles can place very different fractions of the same total mass there.

The implementation in \texttt{EarthBoundAxionHalo} models axions as nonrelativistic bound eigenstates in Earth's gravitational potential. The potential uses a PREM-like interior density profile and a point-mass exterior, with the energy reference set at infinity. For each axion frequency, the code solves the radial Schr\"odinger problem for selected angular-momentum channels and radial states. It then normalizes a state's density to a mass limit in the shell from $12{,}300$ to $384{,}000\,\mathrm{km}$, approximately spanning the outer Earth region to the Moon's orbital distance. If $u(r)$ is the radial wavefunction, the shell probability is $I=\int_{r_1}^{r_2}|u(r)|^2\,dr$; the inferred mass scale is $M_{\rm shell}/I$. The code uses this scale to calculate enclosed mass and the angularly averaged density, including at Earth's surface.

In the example, the shell mass limit is $4\times10^{-9}M_\oplus$. Each eigenstate is treated as an independent profile hypothesis; the code does not add their masses into a simultaneous population. It also applies a constructor-level total-mass cap (normally $0.01M_\oplus$) and an Earth-enclosed-mass cap (normally $0.5M_\oplus$). Thus an inference may be capped below the naive shell rescaling, and the $0.0038M_\oplus$ benchmark is compatible with the default total cap but is not implied by it. A tiny shell probability can formally demand a very large total mass; such cases are sensitive to grid resolution and may violate the neglected-self-gravity approximation. The outer computational radius and radial resolution must therefore be checked for convergence before interpreting the result physically.

The most direct summary figure should show mean DM density at Earth's radius versus axion frequency. Plot frequency on a logarithmic horizontal axis and $\rho_\oplus/\rho_{\rm SHM}$ on a logarithmic vertical axis, with separate curves for the selected eigenstates; here $\rho_{\rm SHM}=0.3\,\mathrm{GeV}/c^2\,\mathrm{cm}^{-3}$. A right axis can report density in $\mathrm{g\,cm^{-3}}$, and a top axis can translate frequency to axion mass using $m_ac^2=h\nu_a$. Such a plot presents the code's shell-normalized, state-by-state predictions across frequency. It does not directly measure total bound mass, and it should be accompanied by the shell limit, applied caps, and convergence range. The defensible conclusion is therefore conditional: a $0.38\%$ Earth-mass halo is a useful benchmark to test, while the allowed surface density and total mass depend on the assumed state, frequency, constraints, and numerical stability.

\end{document}
